% Part: lambda-calculus
% Chapter: lambda-definability
% Section: introduction

\documentclass[../../../include/open-logic-section]{subfiles}

\begin{document}

\olfileid{lam}{ldf}{int}
\olsection{ভূমিকা}

প্রথম দেখায় ল্যাম্বডা ক্যালকুলাস নিছকই এমন অভিব্যক্তির এক অতি বিমূর্ত
ক্যালকুলাস, যেগুলি অপেক্ষক এবং অন্য অপেক্ষকের উপর সেগুলির প্রয়োগকে
উপস্থাপন করে। ল্যাম্বডা ক্যালকুলাসের সংকেতবিন্যাসে এমন কিছু নেই যা বলে
যে এগুলি বিশেষ কোনো ধরনের বস্তুর অপেক্ষক; বিশেষত, এই সংকেতবিন্যাসে
স্বাভাবিক সংখ্যার কোনো উল্লেখই নেই। এর মৌলিক ক্রিয়া---প্রয়োগ ও
ল্যাম্বডা-বিমূর্তন---যেকোনো অপেক্ষকের ক্ষেত্রেই প্রযোজ্য, কেবল স্বাভাবিক
সংখ্যার উপর অপেক্ষকের ক্ষেত্রে নয়।

তবু সামান্য কৌশলে ল্যাম্বডা ক্যালকুলাসে পাটিগাণিতিক অপেক্ষক, অর্থাৎ
স্বাভাবিক সংখ্যার উপর অপেক্ষক, সংজ্ঞায়িত করা যায়। এর জন্য প্রতিটি
স্বাভাবিক সংখ্যা~$n \in \Nat$-এর ক্ষেত্রে আমরা একটি বিশেষ
$\lambd$-পদ~$\num n$ সংজ্ঞায়িত করি, যা $n$-এর \emph{চার্চ সংখ্যাপদ}।
(চার্চ সংখ্যাপদগুলির নাম অ্যালোন্‌জো চার্চের নামে।)

\begin{defn}
  যদি $n \in \Nat$, তবে সংশ্লিষ্ট \emph{চার্চ সংখ্যাপদ} $\num{n}$
  সংখ্যা~$n$-কে উপস্থাপন করে:
  \[
    \num{n} \ident \lambd[fx][f^n(x)]
  \]
  এখানে $f^n(x)$ বলতে $x$-এর উপর $f$-কে $n$ বার প্রয়োগের ফল বোঝায়।
  যেমন, $\num{0}$ হলো $\lambd[fx][x]$, এবং $\num{3}$ হলো
  $\lambd[fx][f(f(f\,x))]$।
\end{defn}

চার্চ সংখ্যাপদ $\num n$ এমন একটি অপেক্ষককে উপস্থাপনকারী ল্যাম্বডা পদ
হিসেবে সাংকেতিত, যে অপেক্ষকটি $f$ ও~$x$ এই দুটি নিবেশ গ্রহণ করে এবং
$f^n(x)$ ফেরত দেয়। চার্চ সংখ্যাপদগুলি স্পষ্টতই স্বাভাবিক রূপে আছে।

ল্যাম্বডা ক্যালকুলাসে স্বাভাবিক সংখ্যার কোনো উপস্থাপন তখনই কাজে লাগে,
যদি তা দিয়ে গণনা করা যায়। ল্যাম্বডা ক্যালকুলাসে চার্চ সংখ্যাপদ দিয়ে
গণনা করার অর্থ হলো এমন কোনো চার্চ সংখ্যাপদের উপর একটি $\lambd$-পদ~$F$
প্রয়োগ করা এবং সম্মিলিত পদ~$F\, \num n$-কে একটি স্বাভাবিক রূপে হ্রাস
করা। যদি এটি সবসময় একটি স্বাভাবিক রূপে হ্রাস পায় এবং সেই স্বাভাবিক রূপ
সবসময় কোনো চার্চ সংখ্যাপদ~$\num m$ হয়, তবে গণনার ফলকে আমরা সংখ্যা~$m$
হিসেবে ভাবতে পারি। তখন~$F$-কে $f\colon \Nat \to \Nat$ অপেক্ষকটি
সংজ্ঞায়িত করছে বলে ভাবতে পারি; অর্থাৎ, সেই অপেক্ষক যার ক্ষেত্রে
$f(n) = m$ যদি এবং কেবল যদি $F\, \num n \red \num m$। চার্চ--রসার ধর্মের
কারণে স্বাভাবিক রূপের অস্তিত্ব থাকলে তা একটিই। তাই
$F\, \num n \red \num m$ হলে, স্বাভাবিক রূপে থাকা অন্য কোনো পদে,
বিশেষত অন্য কোনো চার্চ সংখ্যাপদে, $F \, \num n$ হ্রাস পেতে পারে না।

বিপরীতভাবে, কোনো অপেক্ষক $f\colon \Nat \to \Nat$ দেওয়া থাকলে আমরা
জিজ্ঞাসা করতে পারি, এইভাবে~$f$-কে সংজ্ঞায়িত করে এমন কোনো পদ~$F$ আছে
কি না। থাকলে বলি, $F$ পদটি~$f$-কে \emph{!!{lambda define}s}, এবং~$f$
হলো !!{lambda definable}। এই ধারণাকে বহুস্থানীয় ও আংশিক অপেক্ষকের
ক্ষেত্রেও সাধারণীকরণ করা যায়।

\begin{defn}
ধরা যাক $f\colon \Nat^k \to \Nat$। বলি, কোনো ল্যাম্বডা পদ~$F$
$f$-কে \emph{!!{lambda define}s}, যদি সব
$n_0$, \dots,~$n_{k-1}$-এর জন্য
\[
F \, \num{n_0} \, \num{n_1} \dots \num{n_{k-1}} \red \num{f(n_0, n_1, \dots,
  n_{k-1})}
\]
যখন $f(n_0, \dots, n_{k-1})$ সংজ্ঞায়িত, আর অন্যথায়
$F \, \num{n_0} \, \num{n_1} \dots \num{n_{k-1}}$-এর কোনো স্বাভাবিক রূপ
নেই।
\end{defn}

খুব সরল উদাহরণ হলো ধ্রুবক অপেক্ষকগুলি। $C_k \ident
\lambd[x][\num{k}]$ পদটি $c_k\colon \Nat \to \Nat$ অপেক্ষককে
!!{lambda define}s, যেখানে $c_k(n) = k$।
% BN-SRC-305 TARGET-CORRECTION-BEGIN
\par\smallskip\noindent\textnormal{\small\textbf{উৎস-সংশোধন (BN-SRC-305):} হিমায়িত বাক্যটি অপেক্ষকটির নাম প্রথমে~$c_k$ দিলেও তার নির্ধারক সমীকরণে ভুল করে $c(n)=k$ লেখে। বাংলা পাঠে একই নামযুক্ত অপেক্ষকের অধোলিখন পুনরুদ্ধার করে $c_k(n)=k$ করা হয়েছে।} % BN-SRC-305 TARGET-CORRECTION-END
যেকোনো~$n$-এর জন্য $C_k \, \num n \ident
(\lambd[x][\num{k}])\num n \redone \num{k}$। অভেদী অপেক্ষকটি
$\lambd[x][x]$ দ্বারা !!{lambda defined}। অবশ্যই আরও জটিল অপেক্ষক
সংজ্ঞায়িত করা কঠিন এবং প্রায়ই অনেক কৌশল লাগে। তাই প্রতিটি গণনসাধ্য
অপেক্ষক যে !!{lambda definable}, তা হয়তো বিস্ময়কর। বিপরীতটিও সত্য:
কোনো অপেক্ষক !!{lambda definable} হলে সেটি গণনসাধ্য।

\end{document}
